% GENERATED FILE. Edit canonical lessons, then run npm run generate:books.
% canonical-source-hash: fnv1a64:9e06901e65c8b33b
% canonical-lessons: FR-C341-malentendu, FR-C341-clarifier, FR-C341-preciser, FR-C341-precision, FR-C341-eclaircir, FR-R341-revision-334-341

\chapter{Je m'explique}
\label{ch:fr-je-m-explique}
\hlchaptermodality{341}

\begin{chapteropening}
\textbf{By the end of this chapter:} \emph{I can clear up a misunderstanding: clarify, specify, give more details and clear a question up.}

Complete the last lesson of chapter 341: I can clear up a misunderstanding: clarify, specify, give more details and clear a question up.
\end{chapteropening}

\section[le malentendu]{le malentendu — the misunderstanding}

\label{lesson:FR-C341-malentendu}



\begin{quote}
Before the new one: say the French for likely, then the French for unlikely.
\end{quote}

\subsection*{You'll want to know: le malentendu}
\textbf{le malentendu} — \textquotedblleft{}the misunderstanding\textquotedblright{}.

\begin{quote}
Excuse-moi, c'est un malentendu : je parlais de l'autre réunion.
\end{quote}

\begin{quote}
\emph{Sorry, it is a misunderstanding: I was talking about the other meeting.}
\end{quote}

From \emph{mal} + \emph{entendu}: badly heard.

\begin{cousinweb}
\emph{malentendu} is built from \emph{mal} \textquotedblleft{}badly\textquotedblright{} and \emph{entendu}, the past participle of \emph{entendre}, which comes from Latin \emph{intendere} \textquotedblleft{}to stretch toward, to turn the mind to\textquotedblright{} (Lewis \& Short).
\end{cousinweb}

\begin{grammarlens}[title={What you've built}]
One more word for telling, explaining and describing at B1.
\end{grammarlens}

\subsection*{Your turn}
\begin{itemize}
\raggedright
  \item \emph{Say it:} \emph{le malentendu}
  \item \emph{Say it:} \emph{le malentendu}, once more
  \item \emph{Say it:} a sentence of your own with \emph{le malentendu}
  \item \emph{Recall:} say the French for likely, then the French for unlikely, then say \emph{le malentendu} again
\end{itemize}

\begin{tcolorbox}[breakable,colback=teal!4,colframe=teal!35!black,title={Before you move on}]
What does le malentendu mean? (\textquotedblleft{}The misunderstanding\textquotedblright{}.) Say it once more.
\end{tcolorbox}
\section[clarifier]{clarifier — to clarify}

\label{lesson:FR-C341-clarifier}



\begin{quote}
Before the new one: say the French for unlikely, then the French for the misunderstanding.
\end{quote}

\subsection*{You'll want to know: clarifier}
\textbf{clarifier} — \textquotedblleft{}to clarify\textquotedblright{}.

\begin{quote}
Je voudrais clarifier un point avant de commencer.
\end{quote}

\begin{quote}
\emph{I would like to clarify one point before starting.}
\end{quote}

\emph{Clarifier la situation} — to make things clear.

\begin{cousinweb}
\emph{clarifier} is a learned borrowing from Late Latin \emph{clārificāre} (EtymDB), which Lewis \& Short gloss \textquotedblleft{}to make illustrious, to make famous\textquotedblright{}; it is built on \emph{clārus} \textquotedblleft{}clear, bright\textquotedblright{}, and French took it in the sense to make clear.
\end{cousinweb}

\begin{grammarlens}[title={What you've built}]
One more word for telling, explaining and describing at B1.
\end{grammarlens}

\subsection*{Your turn}
\begin{itemize}
\raggedright
  \item \emph{Say it:} \emph{clarifier}
  \item \emph{Say it:} \emph{clarifier}, once more
  \item \emph{Say it:} a sentence of your own with \emph{clarifier}
  \item \emph{Recall:} say the French for unlikely, then the French for the misunderstanding, then say \emph{clarifier} again
\end{itemize}

\begin{tcolorbox}[breakable,colback=teal!4,colframe=teal!35!black,title={Before you move on}]
What does clarifier mean? (\textquotedblleft{}To clarify\textquotedblright{}.) Say it once more.
\end{tcolorbox}
\section[préciser]{préciser — to specify, to make precise}

\label{lesson:FR-C341-preciser}



\begin{quote}
Before the new one: say the French for the misunderstanding, then the French for to clarify.
\end{quote}

\subsection*{You'll want to know: préciser}
\textbf{préciser} — \textquotedblleft{}to specify, to make precise\textquotedblright{}.

\begin{quote}
Il a précisé l'heure de la réunion dans son message.
\end{quote}

\begin{quote}
\emph{He specified the time of the meeting in his message.}
\end{quote}

\emph{Préciser que} — to point out exactly. From \emph{précis}, precise.

\begin{cousinweb}
\emph{préciser} is built on \emph{précis} \textquotedblleft{}precise\textquotedblright{} (EtymDB), from Latin \emph{praecīsus} \textquotedblleft{}cut off, abridged\textquotedblright{}, the participle of \emph{praecīdere} \textquotedblleft{}to cut off in front\textquotedblright{} (Lewis \& Short).
\end{cousinweb}

\begin{grammarlens}[title={What you've built}]
One more word for telling, explaining and describing at B1.
\end{grammarlens}

\subsection*{Your turn}
\begin{itemize}
\raggedright
  \item \emph{Say it:} \emph{préciser}
  \item \emph{Say it:} \emph{préciser}, once more
  \item \emph{Say it:} a sentence of your own with \emph{préciser}
  \item \emph{Recall:} say the French for the misunderstanding, then the French for to clarify, then say \emph{préciser} again
\end{itemize}

\begin{tcolorbox}[breakable,colback=teal!4,colframe=teal!35!black,title={Before you move on}]
What does préciser mean? (\textquotedblleft{}To specify, to make precise\textquotedblright{}.) Say it once more.
\end{tcolorbox}
\section[la précision]{la précision — the precision; an extra detail}

\label{lesson:FR-C341-precision}



\begin{quote}
Before the new one: say the French for to clarify, then the French for to specify.
\end{quote}

\subsection*{You'll want to know: la précision}
\textbf{la précision} — \textquotedblleft{}the precision; an extra detail\textquotedblright{}.

\begin{quote}
Merci pour ces précisions.
\end{quote}

\begin{quote}
\emph{Thanks for these extra details.}
\end{quote}

In the plural, \emph{des précisions} are extra details: \emph{donner des précisions}.

\begin{cousinweb}
\emph{précision} is a learned borrowing from Latin \emph{praecīsiō} \textquotedblleft{}a cutting off\textquotedblright{} (EtymDB; Lewis \& Short), from \emph{praecīdere} \textquotedblleft{}to cut off in front\textquotedblright{}.
\end{cousinweb}

\begin{grammarlens}[title={What you've built}]
One more word for telling, explaining and describing at B1.
\end{grammarlens}

\subsection*{Your turn}
\begin{itemize}
\raggedright
  \item \emph{Say it:} \emph{la précision}
  \item \emph{Say it:} \emph{la précision}, once more
  \item \emph{Say it:} a sentence of your own with \emph{la précision}
  \item \emph{Recall:} say the French for to clarify, then the French for to specify, then say \emph{la précision} again
\end{itemize}

\begin{tcolorbox}[breakable,colback=teal!4,colframe=teal!35!black,title={Before you move on}]
What does la précision mean? (\textquotedblleft{}The precision; an extra detail\textquotedblright{}.) Say it once more.
\end{tcolorbox}
\section[éclaircir]{éclaircir — to clear up, to make clearer}

\label{lesson:FR-C341-eclaircir}



\begin{quote}
Before the new one: say the French for to specify, then the French for the precision.
\end{quote}

\subsection*{You'll want to know: éclaircir}
\textbf{éclaircir} — \textquotedblleft{}to clear up, to make clearer\textquotedblright{}.

\begin{quote}
Nous devons éclaircir ce problème avant lundi.
\end{quote}

\begin{quote}
\emph{We need to clear up this problem before Monday.}
\end{quote}

A \emph{-ir} verb like \emph{finir}. Of the sky too: \emph{le ciel s'éclaircit}.

\begin{cousinweb}
\emph{éclaircir} is built from \emph{é-} and \emph{clair} \textquotedblleft{}clear\textquotedblright{}, which is inherited from Latin \emph{clārus} \textquotedblleft{}clear, bright\textquotedblright{} (Lewis \& Short).
\end{cousinweb}

\begin{grammarlens}[title={What you've built}]
One more word for telling, explaining and describing at B1.
\end{grammarlens}

\subsection*{Your turn}
\begin{itemize}
\raggedright
  \item \emph{Say it:} \emph{éclaircir}
  \item \emph{Say it:} \emph{éclaircir}, once more
  \item \emph{Say it:} a sentence of your own with \emph{éclaircir}
  \item \emph{Recall:} say the French for to specify, then the French for the precision, then say \emph{éclaircir} again
\end{itemize}

\begin{tcolorbox}[breakable,colback=teal!4,colframe=teal!35!black,title={Before you move on}]
What does éclaircir mean? (\textquotedblleft{}To clear up, to make clearer\textquotedblright{}.) Say it once more.
\end{tcolorbox}
\section[(dialogue)]{Raisons et opinions}

\label{lesson:FR-R341-revision-334-341}



\begin{quote}
Say the French for the point of view and to maintain.
\end{quote}

\subsection*{Your turn}
Say these aloud:

\begin{itemize}
\raggedright
  \item the point of view, to maintain, to claim, to state firmly, convincing
  \item the debate, to debate, the thesis, the line of argument, to persuade
  \item the disagreement, to object, the objection, to retort, to contradict
  \item to admit, the consensus, the compromise, to give in, to concede
  \item to doubt, sceptical, to deny, absurd, debatable
  \item the proof, obvious, the obvious, indisputable, clear
  \item reliable, believable, doubtful, likely, unlikely
  \item the misunderstanding, to clarify, to specify, the precision, to clear up
\end{itemize}

\begin{tcolorbox}[breakable,colback=teal!4,colframe=teal!35!black,title={Before you move on}]
The point of view? (\textbf{le point de vue}.) To maintain? (\textbf{soutenir}.) To claim? (\textbf{prétendre}.) To state firmly? (\textbf{affirmer}.) Convincing? (\textbf{convaincant}.) The debate? (\textbf{le débat}.) To debate? (\textbf{débattre}.) The thesis? (\textbf{la thèse}.) The line of argument? (\textbf{l'argumentation}.) To persuade? (\textbf{persuader}.) The disagreement? (\textbf{le désaccord}.) To object? (\textbf{objecter}.) The objection? (\textbf{l'objection}.) To retort? (\textbf{rétorquer}.) To contradict? (\textbf{contredire}.) To admit? (\textbf{admettre}.) The consensus? (\textbf{le consensus}.) The compromise? (\textbf{le compromis}.) To give in? (\textbf{céder}.) To concede? (\textbf{concéder}.) To doubt? (\textbf{douter}.) Sceptical? (\textbf{sceptique}.) To deny? (\textbf{nier}.) Absurd? (\textbf{absurde}.) Debatable? (\textbf{discutable}.) The proof? (\textbf{la preuve}.) Obvious? (\textbf{évident}.) The obvious? (\textbf{l'évidence}.) Indisputable? (\textbf{incontestable}.) Clear? (\textbf{manifeste}.) Reliable? (\textbf{fiable}.) Believable? (\textbf{crédible}.) Doubtful? (\textbf{douteux}.) Likely? (\textbf{vraisemblable}.) Unlikely? (\textbf{invraisemblable}.) The misunderstanding? (\textbf{le malentendu}.) To clarify? (\textbf{clarifier}.) To specify? (\textbf{préciser}.) The precision? (\textbf{la précision}.) To clear up? (\textbf{éclaircir}.)
\end{tcolorbox}
